\chapter*{Conclusion générale}
\markboth{Conclusion générale}{}
\addcontentsline{toc}{chapter}{Conclusion générale}
\adjustmtc
\thispagestyle{MyStyle}
\begin{onehalfspace}

Ce rapport a présenté la conception et la réalisation d'une plateforme de dématérialisation de la gestion des candidatures, dans le cadre d'un stage de trois mois au Ministère de l'Économie et des Finances, à Rabat, au sein de la Direction des Affaires administratives et générales. Le constat de départ était celui d'une organisation des concours encore largement manuelle : import des listes, affectation des candidats aux salles, édition et envoi des convocations. Cette charge expose à des erreurs, des retards et une faible visibilité d'ensemble.

Le travail a suivi un cycle en cascade, avec conception UML avant le développement, et validation de chaque étape par l'encadrant. Le cahier des charges a précisé les acteurs --- \textbf{gestionnaire} (écriture métier) et \textbf{administrateur} (lecture du métier et gestion des comptes) --- ainsi que les besoins fonctionnels et non fonctionnels. La conception a formalisé les cas d'utilisation, les séquences et le modèle de classes. La réalisation a ensuite mis en œuvre une application web React et un backend de six microservices Spring Boot derrière une passerelle unique, avec une base PostgreSQL par service et une authentification par jeton JWT.

La solution réalisée couvre le périmètre retenu : import des candidats par fichier Excel, gestion des concours et des lieux d'examen, répartition automatique selon la capacité des salles et la proximité géographique, tableau de bord, PDF unitaire de convocation et envoi groupé par courrier électronique. Une convocation n'est assemblée que si l'affectation est complète (centre, établissement, salle et place). Il n'existe ni espace candidat, ni inscription en ligne : l'outil s'adresse aux agents chargés de l'organisation des épreuves.

Ces travaux répondent aux objectifs du stage : disposer d'un système intégré, conçu puis développé, capable d'assister le gestionnaire tout en conservant un contrôle humain. L'application a été réalisée et illustrée dans ce rapport ; elle n'a pas été mise en exploitation par le ministère durant la période du stage.

\end{onehalfspace}

\chapter*{Perspectives}
\markboth{Perspectives}{}
\addcontentsline{toc}{chapter}{Perspectives}
\adjustmtc
\thispagestyle{MyStyle}
\begin{onehalfspace}

Plusieurs évolutions pourraient prolonger ce travail. Sur le plan technique, l'introduction d'un jeton de rafraîchissement renforcerait la gestion des sessions, aujourd'hui limitée à un jeton d'accès. Le caractère non atomique de la répartition pourrait être corrigé, afin qu'un incident en cours d'exécution ne laisse pas d'affectations partielles. L'envoi des convocations, actuellement assuré via Gmail, gagnerait à s'appuyer sur une messagerie institutionnelle.

Sur le plan fonctionnel, un espace de consultation pour le candidat --- téléchargement de sa convocation, sans inscription en ligne --- constituerait une extension naturelle, hors du périmètre actuel. Enfin, une mise en exploitation au sein du ministère supposerait un déploiement maîtrisé, des tests complémentaires et un accompagnement des utilisateurs, dans le respect des règles de sécurité de l'administration.

\end{onehalfspace}
